% 来源：OEBPS/Text/01_36.xhtml（OPF spine 顺序；仅第一册）
\cleardoublepage
\BookEntry[罗伯特·哈格斯特朗《巴菲特的投资组合》]{罗伯特·哈格斯特朗《巴菲特的投资组合》}

\BookCover{images/书单21巴菲特的投资组合.jpg}{[美] 罗伯特·哈格斯特朗著| 江春译| 机械工业出版社2013年版}

\begin{BookSummaryBox}
\textbf{全球第五大上市公司伯克希尔·哈撒韦的企业管理成功方法}
\par
\end{BookSummaryBox}

\subsubsection*{这本书谈什么？}

罗伯特·哈格斯特朗，莱格·梅森集中资产公司副总裁兼总经理，他在这本书中巧妙地解释了众人瞩目的“股票之神”巴菲特是如何通过运用投资组合，而在股市中出奇制胜的。巴菲特集中投资的核心概念主要包括这几个方面：如何找出这些高于平均收益的股票？究竟多少算做“少数几种”？“集中”是什么意思？到底该持股多久，以及为什么要这样做。
\par

作者用浅显易懂的语言说明了不同投资背景的人该如何使用集中投资的策略，为个体投资者提供了成功指南。
\par

\subsubsection*{芒格为什么谈起这本书？}

芒格在谈到为什么没有更多的公司和投资者效仿伯克希尔·哈撒韦公司的问题时，推荐了罗伯特·哈格斯特朗这本著作。因为它清楚地列出了伯克希尔·哈撒韦的成功方法，让每个人都能学习。而关于这些方法，他谈道：
\par

\begin{quotation}
我们的方法让我们取得了成功。看看我们、我们的管理人员和股东有多么开心就知道啦。大多数人应该效仿我们。我们的投资方法并不难学，但表面上看起来很难，因为它跟一般的投资方法不同——它不是常见的方法。我们的管理费用很低，没有季度目标和预算，也没有一套标准的人事体系，我们的投资集中度比一般企业高得多。就这么简单而且合情合理。
\par
\end{quotation}

\subsubsection*{《巴菲特的投资组合》作者罗伯特·哈格斯特朗如是说}

“事实往往是，那些只做几个数量不多的决策的人最终结果，却比那些以忙碌为名作出很多决定的人要好得多。”
\par

“集中投资战略是以经济为基础的模式而不是以价格为基础的模式。投资是否成功和有效，以企业经济效益来衡量，而不是以其股价的短期波动来衡量。”
\par

“在赛马场上或股市上最明智的做法是按兵不动，耐心等待，直到那匹好马出现在有很大获胜机会的位置上才出击。”
\par

“可以毫不夸张地说，整个市场就是被人们的心理因素推来推去。”
\par

“巴菲特和彼得·林奇都说过，学习不过是注意观察你周围发生的事情。”
\par
